%%%%

%%%   Самарский государственный технический университет

%%%   Кафедра прикладной математики и информатики

%%%

%%%   Преамбула для статьи в журнал "Вестник Самарского государственного

%%%   технического университета. Серия: «Физико-математические науки»"

%%%   Ориентирована под MikTEx 2.4 for Win32

%%%

%%%   Хотя сам журнал собирается под GNU/Linux на дистрибутиве TeXLive



\documentclass[11pt,a4paper,twoside]{article}



\usepackage[cp1251]{inputenc}

\usepackage[T2A]{fontenc}

\usepackage[english,russian]{babel}

\usepackage{samgtubib}

\usepackage[dvips]{graphicx}

\usepackage[multiple]{footmisc}



\usepackage{samgtu}

\renewcommand{\VestnikYear}{2009} % Год издания Вестника

\renewcommand{\VestnikNumber}{2\,(19)} % Номер Вестника

\renewcommand{\VestnikMonth}{Ноябрь} % Месяц выхода Вестника

\renewcommand{\VestnikData}{01.11.09} % Дата подписания Вестника в печать

\renewcommand{\VestnikUPL}{26,64} % Усл. печ. л.

\renewcommand{\VestnikUIL}{25,73} % Уч.-изд. л.

\renewcommand{\VestnikRegNumber}{479/09} % Регистрационный номер

\renewcommand{\VestnikOrder}{994} % Номер заказа






%
%\usepackage{pst-barcode}
%\usepackage{auto-pst-pdf}
%%
%%\newcommand{\totop}[1]{\raisebox{6pt}[1pt][2pt]{#1}} 
%%\newcommand{\tottop}[1]{\raisebox{12pt}[1pt][2pt]{#1}} 
%%\newcommand{\totttop}[1]{\raisebox{18pt}[1pt][2pt]{#1}} 
%%\newcommand{\tottttop}[1]{\raisebox{24pt}[1pt][2pt]{#1}} 
%%\newcommand{\totttttop}[1]{\raisebox{30pt}[1pt][2pt]{#1}} 
%\newcommand{\tobot}[1]{\raisebox{-6pt}[1pt][2pt]{#1}} 
%%\newcommand{\tobbot}[1]{\raisebox{-12pt}[1pt][2pt]{#1}} 
%%\newcommand{\tobbbot}[1]{\raisebox{-18pt}[1pt][2pt]{#1}}
%
%\graphicspath{{./00PIC/}}
%
%%\samgtupubl % для генерации pdf, отдаваемого в типографию СамГТУ
%\pdfcrop % обрезка pdf для размещения на math-net.ru
%\draft % На корректуре
%\filllastpage % Последняя страница не заполнена не полностью
%%\briefcomm % Статья является кратким сообщением

\vsgtu{1451}

\FirstPage{667}
\LastPage{679}
%%
%%
%\begin{document}


\renewcommand{\baselinestretch}{0.9}

\hfill \begin{pspicture}(1in,1in)
\psbarcode{http://dx.doi.org/10.14498/vsgtu1451}{}{qrcode}
\end{pspicture}
\vspace{-1in}


%%%%%%%%%%%%%%%%%%%%%%%%%%%%%%%%%%%%%%%%%%%%%%%%%%%%%%%%%%%%%%%%%%%%%%%%%%%%%%%%%%%%
%Информация о статье и об авторах на русском и английском языках

% Номер УДК
\udk{517.956.47}%Обработка текста. Подготовка текстов : Языки программирования. Компьютерные языки
\msc{35G16, 35R30}% Коды MSC см. http://www.ams.org/msc/

\rutitle{Обратная задача для двумерного уравнения теплопроводности по отысканию начального распределения}% Полный заголовок
{Обратная задача для двумерного уравнения теплопроводности\dots}% Заголовок, который идёт в верхний колонтитул

\ruauthor{А.~Р.~Зайнуллов}{З\,а\,й\,н\,у\,л\,л\,о\,в~А.~Р.}


\ruaffil{\rubashgusterl.}
\enaffil{\enbashgusterl.}


\ruauthordetails{1}{\fullauthor{\rupid{113825}{Артур Рашитович Зайнуллов}}\regalia{\mem{arturzayn@mail.ru}}\position{аспирант}
\dept{каф. математического анализа}
 }


\enauthordetails{1}{\fullauthor{\enpid{113825}{Artur R. Zaynullov}}\regalia{\mem{arturzayn@mail.ru}}\position{Postgraduate Student}
\dept{Dept. of Mathematical Analysis}}




\enkeywords{heat equation, first initial-boundary value problem, inverse problem, spectral method, uniqueness, existence, integral equation}




\rukeywords{уравнение теплопроводности, первая начально"=граничная задача, обратная задача, спектральный метод, единственность, существование, интегральное уравнение}

\ruabstract{На основании формулы решения первой начально"=граничной задачи для неоднородного двумерного уравнения теплопроводности изучена обратная задача по отысканию начального распределения. В явном виде строится решение прямой начально"=граничной задачи. Единственность решения прямой начально"=граничной задачи доказана на основании свойства полноты системы собственных функций соответствующей однородной задачи Дирихле для оператора Лапласа. Доказана теорема существования решения прямой начально"=граничной задачи. На основе решения этой задачи исследуется обратная задача, установлен критерий единственности решения обратной задачи. Существование решения обратной задачи эквивалентно сведено к разрешимости интегрального уравнения Фредгольма первого рода.}




\entitle{An inverse problem for two-dimensional equations of finding the thermal conductivity of the initial distribution}


\enauthor{A.~R.~Zaynullov}{Z\,a\,y\,n\,u\,l\,l\,o\,v~A.~R.}

\enabstract{The inverse problem of finding the initial distribution has been studied on the basis of formulas for the solution of the first initial-boundary value problem for the inhomogeneous two-dimentional heat equation. The uniqueness of the solution of the direct initial-boundary value problem has proved with the completeness of the eigenfunctions of the corresponding homogeneous Dirichlet problem for the Laplace operator. The existence theorem for solving direct initial boundary value problem has been proved. Inverse problem has been investigated on the basis of the solution of direct problem, a criterion for the uniqueness of the inverse problem of finding the initial distribution has been proved. The existence of the inverse problem solution has been equivalently reduced to Fredholm integral equation of the first kind.}



\date{10/X/2015} % Дата поступления статьи в редакцию.
\revisiondate{12/XI/2015} % В окончательном виде
\accepteddate{19/XI/2015}





\makerutitle

%%%%%%%%%%%%%%%%%%%%%%%%%%%%%%%%%%%%%%%%%%%%%%%%%%%%%%%%%%%%%%%%%%%%%%%%%%%%%%%%%%%%%
%%Информация о статье и об авторах на русском и английском языках
%
%% Номер УДК
%
%
%
%\institute{Стерлитамакский филиал Башкирского государственного университета, 453103,  Республика Башкортостан, г.~Стерлитамак, пр.~Ленина,~49.}
%% Если несколько, то так  \institute{ Организация 1 \and Организация 2  \and Организация 3}
%
% %E-mail's авторов
%
%\otherinf{}
%
%\keywords{
%}
%
%\workabstract{}
%
%\engtitle{INVERSE PROBLEM OF FINDING THE THERMAL CONDUCTIVITY OF THE INITIAL DISTRIBUTION FOR TWO--DIMENSIONAL HEAT EQUATION}
%
%\engauthor
%
%
%\enginstitute{Sterlitamak Branch of Bashkir State University, 49, Prospect Lenina st., Sterlitamak, Republic of Bashkortostan, 453103}% Место работы каждого из авторов на англ. языке
%% Если несколько, то так  \enginstitute{ Организация 1 \and Организация 2  \and Организация 3}
%
%\otherenginfm{
%
%\engabstract{\language=0%Переключение на англ. язык
%The inverse problem of finding the initial distribution has been studied on the basis of formulas for the solution of the first initial-boundary value problem for the inhomogeneous two-dimentional heat equation. The uniqueness of the solution of the direct initial-boundary value problem has proved  with the completeness of the eigenfunctions of the corresponding homogeneous Dirichlet problem for the Laplace operator. The existence theorem for solving direct initial boundary value problem has been proved. Inverse problem has been investigated on the basis of the solution of direct problem, a criterion for the uniqueness of the inverse problem of finding the initial distribution has been proved. The existence of the inverse problem solution has been equivalently reduced to Fredholm integral equation of the first kind.}
%
%\engkeywords{ heat equation, first initial--boundary value problem, inverse problem, spectral method, uniqueness, existence, integral equation}
%
%\date{ / / } % Дата поступления статьи в редакцию.
%\revisiondate{ / / } % В окончательном виде
%
%%%%%%%%%%%%%%%%%%%%%%%%%%%%%%%%%%%%%%%%%%%%%%%%%%%%%%%%%%%%%%%%%%%%%%%%%%%%%%%%%%%%%
%
%
%\maketitle



{\bf 1. Введение}

Рассмотрим уравнение теплопроводности
\begin{equation}\label{a1}
Lu\equiv u_t-a^2(u_{xx}+u_{yy})=F(x,y,t)
\end{equation}
в параллелепипеде
$$D=\Pi\times(0, T), \quad  \Pi=\{(x,y) \mid \ 0<x<l, 0<y<q\}$$
и следующие задачи для него.

\smallskip

{\small\sc Первая начально"=граничная задача.}
\it  Найти в области $D$ функцию \\ $u(x,y,t),$ удовлетворяющую следующим условиям\/{\rm:}
\begin{gather}\label{a2}
u(x,y,t)\in C(\overline{D})\cap C_{x,y,t}^{2, 2, 1}(D);
\\
\label{a4}
Lu\equiv F(x,y,t),\quad  (x,y,t)\in D;
\\
\label{a4}
u(0, y, t)=u(l,y,t)=0,\quad  0\leqslant y\leqslant q,\quad 0\leqslant t \leqslant T,
\\
\label{a41}
u(x, 0, t)=u(x,q,t)=0,\quad  0\leqslant x\leqslant l,\quad 0\leqslant t \leqslant T,
\\
\label{a5}
u(x,y, 0)=\varphi(x,y),\quad 0\leqslant x\leqslant l,\quad  0\leqslant y\leqslant q,
\end{gather}
где $\varphi(x,y)$ и $F(x,y,t)$ "--- заданные функции\rm.


\smallskip

{\small\sc Обратная задача.} \it Найти функции
$u(x,y,t)$ и $\varphi(x,y),$ удовлетворяющие условиям \eqref{a2}--\eqref{a5} и$,$ кроме того$,$ дополнительному условию
\begin{equation}\label{a6}
u(x_0,y_0, t)=h(t),\quad  0<t_0\leqslant t \leqslant t_1\leqslant T,
\end{equation}
где $(x_0,y_0)$ "--- заданная фиксированная точка прямоугольника $\Pi;$ $t_0$ и $t_1$ "--- заданные действительные числа$;$ $F(x,y,t),$ $h(t)$ "--- заданные функции\rm.

\smallskip


Отметим, что указанная обратная задача изучена в \cite[с.~119]{zayn:1} для однородного одномерного уравнения теплопроводности с~однородными граничными условиями второго рода и доказана теорема её единственности при $x_0=0$ и $x_0=l/\pi$.

В данной работе в явном виде строится решение прямой начально"=граничной задачи \eqref{a2}--\eqref{a5} % (2) -- (6)
и на основе решения этой задачи исследуется обратная задача \eqref{a2}--\eqref{a6}. % (2) -- (7). 
В~известных книгах по уравнениям математической физики  \mbox{[\citen{zayn:2}--\citen{zayn:6}]} единственность решения начально"=граничных задач доказывается методом интегралов энергии. 
Здесь, следуя \cite[c.~111]{zayn:7}, \cite{zayn:8},   единственность решения задачи \eqref{a2}--\eqref{a5} доказана на основании свойства полноты системы собственных функций соответствующей однородной задачи Дирихле для оператора Лапласа. 
Также указаны достаточные условия относительно функций $\varphi(x,y)$ и $F(x,y,t)$, при которых решение поставленной задачи существует и оно определяется в виде суммы двойного ряда Фурье, так как в указанных книгах не приводятся обоснование сходимости ряда Фурье в~классе функций \eqref{a2}.

На основании этих результатов установлен критерий единственности решения задачи \eqref{a2}--\eqref{a6} для любой точки $(x_0,y_0)$ прямоугольника $\Pi$ при $l=q$.

\smallskip

{\bf 2. Единственность и существование прямой задачи}

\begin{theorem}[1] Если существует решение задачи \eqref{a2}--\eqref{a5}$,$ удовлетворяющее условиям
\begin{gather}
\label{b7}
\lim_{x\rightarrow0+}u_x\sin\frac{\pi mx}{l}=\lim_{x\rightarrow l-}u_x\sin\frac{\pi mx}{l}=0,
\\
\label{b8}
\lim_{y\rightarrow0+}u_y\sin\frac{\pi ny}{l}=\lim_{y\rightarrow l-}u_y\sin\frac{\pi ny}{l}=0,
\end{gather}
то оно единственно$.$ \end{theorem}

\smallskip

\begin{proof}
Пусть $\mu_{mn}^2$ "--- собственные значения и $X_{mn}(x,y)$ "--- соответствующие им собственные функции спектральной задачи:
\begin{gather}
\label{a71}v_{xx}+v_{yy}+\mu^2v=0,\quad  (x,y)\in\Pi,
\\
\label{a72}
v\bigl|_{x=0}=v\bigl|_{x=l}=0,\quad  v\bigl|_{y=0}=v\bigl|_{y=q}=0,
\end{gather}
которые определяются по формулам
$$ X_{mn}(x,y)=\frac{2}{l}\sin\mu_mx\sin\mu_ny,$$
$$  \mu_{mn}^2=\mu_m^2+\mu_n^2, \mu_m=\frac{\pi m}{l},\quad  \mu_n=\frac{\pi n}{l},\quad  m, n=1, 2, \ldots.$$

Наряду с этой системой рассмотрим  систему
\begin{equation}\label{a73}Z_{mn}(x,y)=\frac{1}{2}\left[X_{mn}(x,y)+X_{nm}(x,y)\right],\end{equation}
которая обладает свойством симметрии
$$Z_{mn}(x,y)=Z_{nm}(x,y).$$

Система \eqref{a73} является решением спектральной задачи \eqref{a71} и \eqref{a72}, обладает свойствами ортогональности и~полноты в пространстве $L_2(\Pi)$.

Рассмотрим интеграл
\begin{equation}\label{a7}
v_{mn}^{\varepsilon,\delta}(t)=\int_{\varepsilon}^{l-\varepsilon}\int_{\delta}^{l-\delta} u(x,y,t)Z_{mn}(x,y)dxdy,
\end{equation}
где $\varepsilon$, $\delta$ "--- достаточно малые положительные числа.

Дифференцируя равенство \eqref{a7} при $t\in(0, T)$ и учитывая уравнение \eqref{a1}, 
получим
\begin{multline}
\label{a8}
v'_{mn}(t)=\int_{\varepsilon}^{l-\varepsilon}\int_{\delta}^{l-\delta} u_t(x,y,t)Z_{mn}(x,y)dxdy =
\\ 
=\int_{\varepsilon}^{l-\varepsilon}\int_{\delta}^{l-\delta} (a^2\Delta u+ F(x,y,t))Z_{mn}(x,y)dxdy =
\\
=a^2\int_{\varepsilon}^{l-\varepsilon}\int_{\delta}^{l-\delta}u_{xx}Z_{mn} x,y)dxdy+ a^2\int_{\varepsilon}^{l-\varepsilon}\int_{\delta}^{l-\delta} u_{yy}Z_{mn}(x,y)dxdy+
\\
+ \int_{\varepsilon}^{l-\varepsilon}\int_{\delta}^{l-\delta}F(x,y,t)Z_{mn}(x,y)dxdy =
\\
=a^2(I_1+I_2)+F_{mn}(t),
\end{multline}
где
\begin{equation}
\label{b9}
F_{mn}=\int_{\varepsilon}^{l-\varepsilon}\int_{\delta}^{l-\delta}F(x,y,t)Z_{mn}(x,y)dxdy.
\end{equation}

С учетом граничных условий \eqref{a4}, \eqref{a41}, $X(0)=X(l)=0$, $Y(0)=Y(q)=0$ проинтегрируем по частям интегралы $I_1$ и $I_2$ в области $\Pi$:
\begin{multline}
I_{1}=\int_{\delta}^{l-\delta}dy\int_{\varepsilon}^{l-\varepsilon} u_{xx}Z_{mn}(x,y)dx= 
\\
=
\int_{\delta}^{l-\delta} dy\biggl(u_x Z_{mn}(x,y)\big|_{\varepsilon}^{l-\varepsilon}-
\int_{\varepsilon}^{l-\varepsilon}u_x\left(Z_{mn}(x,y)\right)'_x dx\biggr)=
\\
=\int_{\delta}^{l-\delta} dy\biggl(u_x Z_{mn}(x,y)\big|_{\varepsilon}^{l-\varepsilon}-u(x,y,t) \left(Z_{mn}(x,y)\right)'\big|_{\varepsilon}^{l-\varepsilon}+
\\
+ \int_{\varepsilon}^{l-\varepsilon}u(x,y,t)\left(Z_{mn}(x,y)\right)''_{xx} dx\biggr),
\end{multline}

\begin{multline}
I_{2}=\int_{\varepsilon}^{l-\varepsilon}dx\int_{\delta}^{l-\delta} u_{yy} Z_{mn}(x,y)dy= 
\\
=
\int_{\varepsilon}^{l-\varepsilon} dx\biggl(u_y Z_{mn}(x,y)\big|_{\delta}^{l-\delta}-
\int_{\delta}^{l-\delta} u_y\left(Z_{mn}(x,y)\right)'_{y} dy\biggr)  =
\\
=\int_{\varepsilon}^{l-\varepsilon} dx\biggl(u_y Z_{mn}(x,y)\big|_{\delta}^{l-\delta}-
u(x,y,t)\left(Z_{mn}(x,y)\right)'_{y} \big|_{\delta}^{l-\delta}+
\\
+
\int_{\delta}^{l-\delta} u(x,y,t)\left(Z_{mn}(x,y)\right)''_{yy} dy\biggr). 
\end{multline}

Из \eqref{a8}  с учетом полученных значений интегралов $I_{1}$ и $I_{2}$ при $\varepsilon\rightarrow 0$, $\delta\rightarrow 0$, следует
\begin{multline}
v'(t)=a^2\iint_\Pi u(x,y,t)\left(\left(Z_{mn}(x,y)\right)''_{xx}+\left(Z_{mn}(x,y)\right)''_{yy}\right)dxdy+F_{mn}(t)=
\\
=-\mu_{mn}^2a^2\iint_\Pi u(x,y,t)Z_{mn}(x,y)dxdy+F_{mn}(t)=
\\
=-\mu_{mn}^2a^2v_{mn}(t)+F_{mn}(t),
\end{multline}
откуда
\begin{equation}\label{a9}v'_{mn}(t)+\mu_{mn}^2a^2v_{mn}(t)=F_{mn}(t).\end{equation}

Общее решение уравнения \eqref{a9} определяется по формуле
\begin{equation}
\label{b10}
v_{mn}(t)=\int_0^tF_{mn}(s)e^{-\mu_{mn}^2a^2(t-s)}ds+C_{mn}e^{-\mu_{mn}^2a^2t},
\end{equation}
где $C_{mn}$ "--- произвольные постоянные. 
Для определения их воспользуемся граничным условием \eqref{a5} и формулой \eqref{a7}:
\begin{equation}
\label{b11}
v_{mn}(0)=\iint_\Pi u(x,y,0)Z_{mn}(x,y)dxdy=\iint_\Pi \varphi(x,y)Z_{mn}(x,y)dxdy=\varphi_{mn}.
\end{equation}

Удовлетворяя \eqref{b10} начальному условию \eqref{b11}, найдем $C_{mn}=\varphi_{mn}$ и
\begin{equation}
\label{b12}
v_{mn}(t)=\varphi_{mn}e^{-\mu_{mn}^2a^2t}+\int_0^tF_{mn}(s)e^{-\mu_{mn}^2a^2(t-s)}ds.
\end{equation}

Отсюда следует единственность решения задачи \eqref{a2}--\eqref{a5}, так как если положить $\varphi(x,y)\equiv0$, 
 $F(x,y,t)\equiv0$, то $\varphi_{mn}\equiv0$, $F_{mn}(s)\equiv0$, и из \eqref{b12} получим $v_{mn}(t)=0$, что на основании \eqref{a7} равносильно равенству
$$\iint_\Pi u(x,y,t)Z_{mn}(x,y)dxdy=0,$$
и в силу полноты системы $Z_{mn}(x,y)$ в пространстве $L_2(\Pi)$ функция ${u(x,y,t){=}0}$ почти всюду в $\Pi$ и при любом $t\in[0, T]$. 
Тем самым единственность решения задачи \eqref{a2}--\eqref{a5} доказана.
\end{proof}

\smallskip

\hypertarget{zay:th2}{}

\begin{theorem}[2] 
Если $\varphi(x,y)\in C^4(\overline{\Pi}),$ $ F(x,y,t)\in C^{2,2}_{x,y}(\overline{\Pi})\cap C_t(0,T)$ и
\begin{gather}
\label{a10}
\varphi(0, y)=\varphi(l,y)=\varphi(x, 0)=\varphi(x,l)=0,\\
\varphi'_x(0, y)=\varphi'_x(l,y)=\varphi'_x(x, 0)=\varphi'_x(x,l)=0,\\
\varphi'_y(0, y)=\varphi'_y(l,y)=\varphi'_y(x, 0)=\varphi'_y(x,l)=0,\\
\label{a11}
F(0, y,t)=F(l,y,t)=F(x, 0, t)=F(x,l,t)=0,\\
F'_x(0, y,t)=F'_x(l,y,t)=F'_x(x, 0, t)=F'_x(x,l,t)=0,\\
F'_y(0, y,t)=F'_y(l,y,t)=F'_y(x, 0, t)=F'_y(x,l,t)=0
\end{gather}
для всех $(x,y,t)\in \overline{D},$ то существует решение задачи \eqref{a2}--\eqref{a5} и оно представимо в виде суммы ряда
\begin{equation}
\label{a12}
u(x,y,t)=\sum_{m,n=1}^\infty\left[\varphi_{mn}e^{-\mu_{mn}^2a^2t}+\int_0^tF_{mn}(s)e^{-\mu_{mn}^2a^2(t-s)}ds\right]Z_{mn}(x,y)\end{equation}
с коэффициентами$,$ которые определяются по формулам \eqref{b11} и \eqref{b9}.
\end{theorem}

\smallskip

\begin{proof}
Решение задачи \eqref{a2}--\eqref{a5} выше формально построено в виде суммы ряда \eqref{a12}, где коэффициенты $\varphi_{mn},$ 
$ F_{mn}(t)$ определяются по формулам \eqref{b11} и \eqref{b9}. 
Если ряд \eqref{a12} равномерно сходится на $\overline{D}$, так же как ряды, полученные из него почленным дифференцированием один раз по переменной $t$  и по два раза по $x$, $y$, то его сумма будет удовлетворять в области $D$ условиям \eqref{a2}--\eqref{a5}.

Предварительно проинтегрируем по частям интегралы \eqref{b11}, \eqref{b9} два раза по $x$  и два раза $y$ и, учитывая \eqref{a10}, \eqref{a11}, получим
\begin{multline}
\label{a15}
\varphi_{mn}=\iint_\Pi \varphi(x,y)Z_{mn}(x,y)dxdy= 
\\
=
\frac{1}{l}\int_0^l\sin\mu_nydy\int_0^l\varphi(x,y)\sin\mu_mxdx+ \hspace{3cm}
\\
\hspace{3cm}
+\frac{1}{l}\int_0^l\sin\mu_mydy\int_0^l\varphi(x,y)\sin\mu_nxdx=
\\
=
-\frac{1}{\mu_m^2}\frac{1}{l}\int_0^l\sin\mu_nydy\int_0^l\varphi''_{xx}\sin\mu_mxdx-
\hspace{3cm}
\\
\hspace{3cm}
-\frac{1}{\mu_n^2}\frac{1}{l}\int_0^l\sin\mu_mydy\int_0^l\varphi''_{xx}\sin\mu_nxdx=
\\
= -\frac{1}{\mu_m^2}\frac{1}{l}\int_0^l\sin\mu_mxdx\int_0^l\varphi''_{xx}\sin\mu_nydy-
\hspace{3cm}
\\
\hspace{3cm}
-\frac{1}{\mu_n^2}\frac{1}{l}\int_0^l\sin\mu_nxdx\int_0^l\varphi''_{xx}\sin\mu_mydy=
\\
=
\frac{1}{\mu_m^2\mu_n^2}\frac{1}{l}\int_0^l\sin\mu_mxdx\int_0^l\varphi^{2,2}_{xy}\sin\mu_nydy
+\hspace{3cm}
\\
\hspace{3cm}
+\frac{1}{\mu_m^2\mu_n^2}\frac{1}{l}\int_0^l\sin\mu_nxdx\int_0^l\varphi^{2,2}_{xy}\sin\mu_mydy
=
\\
=
\frac{1}{\mu_m^2\mu_n^2}\frac{1}{l}\iint_\Pi \varphi^{2,2}_{xy}\sin\mu_mx\sin\mu_nydxdy+
\hspace{3cm}
\\
\hspace{3cm}
+\frac{1}{\mu_m^2\mu_n^2}\frac{1}{l}\iint_\Pi \varphi^{2,2}_{xy}\sin\mu_nx\sin\mu_mydxdy=\\
=\frac{1}{\mu_m^2\mu_n^2}\iint_\Pi \varphi^{2,2}_{xy}Z_{mn}(x,y)dxdy
=\frac{l^4}{\pi^4}\frac{\varphi_{mn}^{(4)}}{m^2n^2},
\end{multline}
\begin{multline}
\label{a16}
F_{mn}(t)=\iint_\Pi F(x,y,t)Z_{mn}(x,y)dxdy= 
\\
=
\frac{1}{l}\int_0^l\sin\mu_nydy\int_0^lF(x,y,t)\sin\mu_mxdx+
\hspace{3cm}
\\
\hspace{3cm}
+\frac{1}{l}\int_0^l\sin\mu_mydy\int_0^lF(x,y,t)\sin\mu_nxdx=
\\
=
-\frac{1}{\mu_m^2}\frac{1}{l}\int_0^l\sin\mu_nydy\int_0^lF''_{xx}\sin\mu_mxdx-
\hspace{3cm}
\\
\hspace{3cm}
-\frac{1}{\mu_n^2}\frac{1}{l}\int_0^l\sin\mu_mydy\int_0^lF''_{xx}\sin\mu_nxdx=
\\
=
 -\frac{1}{\mu_m^2}\frac{1}{l}\int_0^l\sin\mu_mxdx\int_0^lF''_{xx}\sin\mu_nydy-
\hspace{3cm} 
 \\
\hspace{3cm} 
 -\frac{1}{\mu_n^2}\frac{1}{l}\int_0^l\sin\mu_nxdx\int_0^lF''_{xx}\sin\mu_mydy
=
\\ 
 =\frac{1}{\mu_m^2\mu_n^2}\frac{1}{l}\int_0^l\sin\mu_mxdx\int_0^lF^{2,2}_{xy}\sin\mu_nydy+
 \hspace{3cm}
\\
\hspace{3cm}
+\frac{1}{\mu_m^2\mu_n^2}\frac{1}{l}\int_0^l\sin\mu_nxdx\int_0^lF^{2,2}_{xy}\sin\mu_mydy=
\\
=
\frac{1}{\mu_m^2\mu_n^2}\frac{1}{l}\iint_\Pi F^{2,2}_{xy}\sin\mu_mx\sin\mu_nydxdy+
\hspace{3cm}
\\
\hspace{3cm}
+\frac{1}{\mu_m^2\mu_n^2}\frac{1}{l}\iint_\Pi F^{2,2}_{xy}\sin\mu_nx\sin\mu_mydxdy=
\\
=
\frac{1}{\mu_m^2\mu_n^2}\iint_\Pi F^{2,2}_{xy}Z_{mn}(x,y)dxdy=\frac{l^4}{\pi^4}\frac{F_{mn}(t)^{(4)}}{m^2n^2}.
\end{multline}

Поскольку $\varphi^{2,2}_{xy}$ и $F^{2,2}_{xy}(t)$ непрерывны в $\overline{\Pi}$, в силу неравенства Бесселя следующие ряды сходятся:
\begin{gather}
\label{a17}
\sum_{m,n=1}^{\infty}|\varphi^{(4)}_{mn}|^2\leqslant\frac{2}{l}\int_0^{l}\int_0^{l}(\varphi^{2,2}_{xy})^2dxdy, 
\\ 
\sum_{m,n=1}^{\infty}|F_{mn}^{(4)}(t)|^2\leqslant\frac{2}{l}\int_0^{l}\int_0^{l}(F^{2,2}_{xy}(t))^2dxdy.
\end{gather}

Подставив \eqref{a15} и \eqref{a16} в ряд \eqref{a12}, получим
\begin{equation}
\label{a18}
u(x,y,t)=\frac{l^4}{\pi^4}\sum_{m,n=1}^\infty\frac{1}{m^2n^2}\biggl[\varphi^{(4)}_{mn}e^{-\mu_{mn}^2a^2t}+\int_0^tF^{(4)}_{mn}(s)e^{-\mu_{mn}^2a^2(t-s)}ds\biggr]
Z_{mn}(x,y).
\end{equation}
Теперь оценим
\begin{multline}
\biggl|\int_0^tF^{(4)}_{mn}(s)e^{-\mu_{mn}^2a^2(t-s)}ds\biggr|=
\max\limits_{0\leqslant s\leqslant T}
|F^{(4)}_{mn}(s)|\int_0^te^{-\mu_{mn}^2a^2(t-s)}ds=
\\
=\frac{\|F^{(4)}_{mn}\|_C(1-e^{-\mu_{mn}^2a^2t_0})}{\mu_{mn}^2a^2}<\frac{l^2}{a^2\pi^2}\frac{\|F^{(4)}_{mn}\|_C}{m^2+n^2}.
\end{multline}
Ряд \eqref{a18} при любых $(x,y,t)\in \overline{D}$ мажорируется сходящимся рядом
\begin{equation}
\label{a19}
\frac{l^4}{\pi^4}\sum_{m,n=1}^{\infty}\frac{1}{m^2n^2}|\varphi^{(4)}_{mn}|+\frac{l^6}{\pi^6a^2}\sum_{m=1}^\infty\frac{1}{m^2+n^2} \sum_{n=1}^\infty\frac{\|F^{(4)}_{mn}\|_C}{m^2n^2}.
\end{equation}
Тогда ряд \eqref{a12} в силу признака Вейерштрасса сходится абсолютно и равномерно в $\overline{D}$. 
Следовательно, функция $u(x,y,t)$ непрерывна в $\overline{D}$ как сумма равномерно сходящегося ряда \eqref{a12}.

Теперь докажем возможность почленного дифференцирования ряда \eqref{a12} по переменным $x$, $y$ два раза и по $t$ один раз в $D$. 
Для этого покажем, что полученные при почленном дифференцировании ряды сходятся абсолютно и равномерно на $\overline{D}_0=\overline{D}\cap\{t\geqslant t_0>0\}$, где $t_0$ "--- достаточно малое положительное число. 
Формально из \eqref{a18} почленным дифференцированием составим следующие ряды:
\begin{multline}
\label{a20}u_{xx}=-\frac{l^4}{\pi^4}\sum_{m,n=1}^\infty\frac{\mu_m^2}{m^2n^2}\biggl(\varphi^{(4)}_{mn}e^{-\mu_{mn}^2a^2t}+
\\
+
\int_0^tF^{(4)}_{mn}(s)e^{-\mu_{mn}^2a^2(t-s)}ds\biggr)Z_{mn}(x,y),
\end{multline}
\begin{multline}
\label{a22}
u_{yy}=-\frac{l^4}{\pi^4}\sum_{m,n=1}^\infty\frac{\mu_n^2}{m^2n^2}\biggl(\varphi^{(4)}_{mn}e^{-\mu_{mn}^2a^2t}+ \\+
\int_0^tF^{(4)}_{mn}(s)e^{-\mu_{mn}^2a^2(t-s)}ds\biggr)Z_{mn}(x,y),
\end{multline}
\begin{multline}\label{a22}u_t=\frac{l^4}{\pi^4}\sum_{m,n=1}^\infty\frac{\mu_{mn}^2a^2}{m^2n^2}\biggl(-\varphi^{(4)}_{mn}e^{-\mu_{mn}^2a^2t}+ 
\\
+\int_0^tF^{(4)}_{mn}(s)e^{-\mu_{mn}^2a^2(t-s)}ds+\frac{F_{mn}^{(4)}(t)}{\mu_{mn}^2a^2}\biggr) Z_{mn}(x,y).
\end{multline}

Ряды \eqref{a20}--\eqref{a22} при любых $(x,y,t)\in \overline{D}_0$ мажорируются соответственно рядами
\begin{gather}
\label{a23}\frac{l^4}{\pi^4}\sum_{m,n=1}^{\infty}\frac{1}{m^2n^2}|\varphi^{(4)}_{mn}|\mu_{m}^2e^{-\mu_{mn}^2a^2t}+\frac{l^5}{\pi^5a^2}\sum_{m=1}^\infty\frac{1}{m(m^2+n^2)}
\sum_{n=1}^{\infty}\frac{\|F^{(4)}_{mn}\|_C}{n^2},
\\
\label{a25}
\frac{l^4}{\pi^4}\sum_{m,n=1}^{\infty}\frac{1}{m^2n^2}|\varphi^{(4)}_{mn}|\mu_{n}^2e^{-\mu_{mn}^2a^2t}+\frac{l^5}{\pi^5a^2}\sum_{n=1}^\infty\frac{1}{n(m^2+n^2)}
\sum_{m=1}^{\infty}\frac{\|F^{(4)}_{mn}\|_C}{m^2},
\\
\label{a25}\frac{l^4a^2}{\pi^4}\sum_{m,n=1}^{\infty}\frac{1}{m^2n^2}|\varphi^{(4)}_{mn}|\mu_{mn}^2e^{-\mu_{mn}^2a^2t}+\frac{2l^4}{\pi^4}
\sum_{m=1}^{\infty}\frac{1}{m^2}\sum_{n=1}^{\infty}\frac{\|F^{(4)}_{mn}\|_C}{n^2}.
\end{gather}

Ряды \eqref{a23}--\eqref{a25} сходятся абсолютно. 
Тогда ряды \eqref{a20}--\eqref{a22} на основании признака Вейерштрасса сходятся абсолютно и равномерно на $D$. 
Следовательно, функции $u_{xx}$, $u_{yy}$, $u_t\in C(D)$ и, подставляя ряды в уравнение \eqref{a12}, убеждаемся в том, что функция $u(x,y,t)$, определяемая рядом \eqref{a12}, является его решением в $D.$ 
\end{proof}

\smallskip

{\bf 3. Критерий единственности решения обратной задачи}

Рассмотрим теперь обратную задачу \eqref{a2}--\eqref{a6}. 
Полагая в формуле \eqref{a12} \linebreak ${x=x_0}$, $y=y_0$ с учётом условия \eqref{a6}, получим относительно неизвестной функции $\varphi(x,y)$ уравнение
\begin{multline}
\label{b6}\sum_{m,n=1}^{\infty}\iint_{\Pi}\varphi(\xi,\eta)Z_{mn}(\xi,\eta)d\xi d\eta e^{-\mu_{mn}^2a^2t}Z_{mn}(x_0,y_0)=
\\
=h(t)-\sum_{m,n=1}^{\infty}F_{mn}(t)Z_{mn}(x_0,y_0)=h_0(t),\quad  t_0\leqslant t\leqslant t_1.
\end{multline}

\smallskip

\hypertarget{zay:th3}{}

\begin{theorem}[3] \label{t22}
Если $Z_{mn}(x_0,y_0)\neq 0$ при всех $m,n\in \mathbb{N},$
то решение интегрального уравнения \eqref{b6} единственно в $L_2(\overline{\Pi})$\rm.
\end{theorem}


\begin{proof}
Следуя \cite[c.~119]{zayn:1}, \cite{zayn:7, zayn:9}, в силу линейности уравнения \eqref{b6} достаточно показать, 
что оно имеет только нулевое решение при $h_0(t) =0$. 
Положив в \eqref{b6} $h(t) = 0$ и $F(x,y,t)\equiv 0$, имеем
\begin{equation}
\label{b7}
\sum_{m,n=1}^{\infty}\int\int_{\Pi}\varphi(\xi,\eta)Z_{mn}(\xi,\eta)d\xi d\eta e^{-\mu_{mn}^2a^2t}Z_{mn}(x_0,y_0)=0.
\end{equation}

Рассмотрим в комплексной полуплоскости $\mathop{\rm Re} z\geqslant\alpha$, где ${\alpha\in(0,t_0)}$ "--- постоянная, функцию комплексной переменной
\begin{equation}\label{b8}
\Phi(z)=\sum_{m,n=1}^{\infty}\varphi_{mn}e^{-\mu_{mn}^2a^2z}Z_{mn}(x_0,y_0).
\end{equation} 
Так как при $\mathop{\rm Re} z\geqslant\alpha$
\begin{equation}\label{b92}
|\varphi_{mn}Z_{mn}(x_0,y_0)e^{-\mu_{mn}^2a^2z}|\leqslant Ce^{-\mu_{mn}^2a^2\alpha},
\end{equation}
где $C=\rm const>0$, то в этой полуплоскости ряд, стоящий в правой части соотношения \eqref{b92}, сходится равномерно. 
Учитывая, что каждый член этого ряда является аналитической функцией при $\mathop{\rm Re} z\geqslant\alpha$, и применяя теорему Вейерштрасса, получаем, что функция $\Phi(z)$ является аналитической при $\mathop{\rm Re} z\geqslant\alpha$. 
Поскольку в силу \eqref{b8} $\Phi(z)=0$ на отрезке $[t_0, t_1]$ действительной оси $t$ из области аналитичности $\Phi(z)$, 
то на основании теоремы единственности для аналитических функций следует $\Phi(z)\equiv 0$ при $\mathop{\rm Re}  z\geqslant \alpha$.
Отсюда следует равенство при всех $t\geqslant t_0$:
\begin{equation}\label{b81}
\sum_{m,n=1}^{\infty}\varphi_{mn}e^{-\mu_{mn}^2a^2t}Z_{mn}(x_0,y_0)=0.
\end{equation}

Умножим равенство \eqref{b81} на $e^{(\mu_{11} a)^2 t}$ и в полученном равенстве, переходя к пределу при $t \rightarrow +\infty$, найдем
\begin{equation}
\label{d1}\varphi_{11}=0.
\end{equation}
С учетом \eqref{d1} умножим равенство \eqref{b81} на $e^{(\mu_{12} a)^2 t}$ и в полученном равенстве, переходя к пределу при $t \rightarrow +\infty$, получим
$$
\varphi_{12}Z_{12}(x_0,y_0)+\varphi_{21}Z_{21}(x_0,y_0)=2\varphi_{12}Z_{12}(x_0,y_0)=0.
$$
Отсюда с учетом того, что $Z_{12}(x_0,y_0)\neq0$, имеем
\begin{equation}
\label{d2}
\varphi_{12}=\varphi_{21}=0.
\end{equation}

Снова учитывая \eqref{d1} и \eqref{d2},  умножая равенство \eqref{b81} на $e^{(\mu_{22} a)^2 t}$ и  переходя к пределу при $t \rightarrow +\infty$ в полученном равенстве, найдем
\begin{equation}
\label{d3}\varphi_{22}=0.
\end{equation}

Затем, учитывая равенства \eqref{d1}--\eqref{d3}, 
умножим равенство \eqref{b81} на $e^{(\mu_{13} a)^2 t}$ и, переходя к пределу при $t\rightarrow \infty$, будем иметь
\begin{equation}
\label{d4}\varphi_{13}=\varphi_{31}=0.
\end{equation}

Затем в силу \eqref{d1}--\eqref{d4} получим
$$
\varphi_{23}=\varphi_{32}=0.
$$

Рассуждая далее аналогично вышеизложенному, получим
\begin{equation}
\label{b100}
\varphi_{mn}=0,\quad  m,n=1, 2,\ldots .
\end{equation}
Из равенства \eqref{b100} в силу полноты системы функций $\{Z_{mn}(x,y)\}_{m,n\geqslant 1}$ в пространстве $L_2(\overline{\Pi})$ следует, что $\varphi(x,y)=0$ почти всюду в $\Pi$. 
Отсюда в силу непрерывности функции $\varphi(x,y)$ на $\overline{\Pi}$ получим, что $\varphi(x,y)\equiv0.$ 
\end{proof}

\smallskip

Тогда из теорем \hyperlink{zay:th3}{3}  и \hyperlink{zay:th2}{2}  следует единственность решения обратной задачи~\eqref{a2}--\eqref{a6}.

\smallskip

Пусть при некоторых $\widetilde{x}_0={x_0}/{l}$, $\widetilde{y}_0={y_0}/{l}$ и $m=p$,  $ n=s\in \mathbb{N}$ нарушено условие теоремы~\hyperlink{zay:th3}{3}:
\begin{equation}
\label{b102}Z_{ps}(x_0,y_0)=\sin\pi p\widetilde{x}_0\sin\pi s\widetilde{y}_0+\sin\pi s\widetilde{x}_0\sin\pi p\widetilde{y}_0=0.\end{equation}
Тогда обратная задача \eqref{a2}--\eqref{a6} при $h(t)=0$ и $F(x,y,t)\equiv 0$ имеет ненулевое решение
\begin{equation}
\label{b112}
u_{ps}(x,y,t)=Z_{ps}(x,y)e^{-(\mu_{ps} a)^2t},\quad  \varphi_{ps}(x,y)=u_{ps}(x,y,0)=Z_{ps}(x,y).
\end{equation}

Из уравнения \eqref{b102} найдем рациональные значения
\begin{equation}
\label{b122}
\widetilde{x}_0={k_1}/{p},\quad   \widetilde{y}_0={k_2}/{p},\quad  k_1 , k_2\in \mathbb{N},\quad  k_1<p, k_2<p,
\end{equation}
или
\begin{equation}
\label{b14}
\widetilde{x}_0={k_3}/{s},\quad   \widetilde{y}_0={k_4}/{s}, \quad  k_3, k_4\in \mathbb{N},\quad  k_3<s, k_4<s,
\end{equation}
при которых нарушается условие теоремы~\hyperlink{zay:th3}{3}, т.е. нарушается единственность решения задачи~\eqref{a2}--\eqref{a6}.

Следовательно, установлен следующий критерий единственности решения обратной задачи~\eqref{a2}--\eqref{a6}.

\smallskip

\begin{theorem}[4] Условия $Z_{mn}(x_0,y_0)\neq0$ при всех $m,n\in\mathbb{N}$ необходимы и достаточны для единственности решения обратной задачи {\em\eqref{a2}--\eqref{a6}}.
\end{theorem}

\smallskip

Рассматриваемая обратная задача может быть сведена к интегральному уравнению Фредгольма первого рода. Действительно, поменяв местами порядок интегрирования и суммирования в левой части уравнения (\ref{b6}), получим интегральное уравнение Фредгольма первого рода
\begin{equation}
\label{b13}
G\varphi\equiv\iint_{\Pi}G(t,\xi,\eta)\varphi(\xi,\eta)d\xi d\eta=h_0(t),\ t_0\leqslant t\leqslant t_1,\end{equation}
с гладким ядром
\begin{equation}
G(t,\xi,\eta)=\sum_{m,n=1}^{\infty} e^{-\mu_{mn}^2a^2t}Z_{mn}(\xi,\eta)Z_{mn}(x_0,y_0)
\end{equation}
в параллелепипеде $\overline{\Pi}\times[t_0,t_1]$. 
Следовательно, интегральный оператор $G$, рассматриваемый действующим из $L_2(\overline{\Pi})$ в $L_2[t_0,t_1]$, вполне непрерывен. 
Задача решения уравнения \eqref{b13} в этой паре пространств некорректна.

\smallskip

{\bf 4. Заключение}

В заключение отметим, что обратные задачи для уравнений параболического типа, как правило, являются некорректными. 
Для их решения обычно привлекается теория некорректных задач, разработанная в работах А.~Н.~Тихонова, М.~М.~Лаврентьева, В.~К.~Иванова и их учеников [\citen{zayn:10}--\citen{zayn:12}]. 
В~данной работе, в отличие от других работ, предлагается спектральный метод для обоснования единственности решения обратной задачи \eqref{a2}--\eqref{a6}. % (2)--(7).



\thanks{\textbf{Благодарности.} 
Работа выполнена при поддержке Российского фонда фундаментальных исследований  (проект \No~14--01--97003-р\_поволжье\_а).}


\thanks{{\bf ORCID}

{\bf Артур Рашитович Зайнуллов:}    \url{http://orcid.org/0000-0003-2657-4101}



}


\newpage

\RusRef


\begin{thebibliography}{99}

\RBibitem{zayn:1}
\by  Денисов~А.~М.
\book Введение в теорию обратных задач
\publaddr М.
\publ МГУ
\yr 1994
\totalpages 208 

\RBibitem{zayn:2}
\by  Тихонов~А.~Н., Самарский~А.~Н.
\book Уравнения математической физики
\publaddr М.
\publ Физматлит
\yr 1966
\totalpages 724 
\zmath{https://zbmath.org/?q=an:0044.09302}

\RBibitem{zayn:3}
\by  Соболев~С.~Л.
\book Уравнения математической физики
\publaddr М.
\publ Наука
\yr 1992
\totalpages 432
\zmath{https://zbmath.org/?q=an:0900.35001}

\RBibitem{zayn:4}
\by Владимиров~В.~С., Жаринов~В.~В.
\book Уравнения математической физики
\publaddr М.
\publ Физматлит
\yr 2000
\totalpages 398 
\zmath{https://zbmath.org/?q=an:0954.35001}

\RBibitem{zayn:5}
\by Кошляков~Н.~С., Глинер~Э.~Б., Смирнов~М.~М.
\book Уравнения в частных производных математической физики
\publaddr М.
\publ Высшая школа
\yr 1970
\totalpages 712 
\zmath{https://zbmath.org/?q=an:0202.10401}

\RBibitem{zayn:6}
\by Годунов~С.~К.
\book Уравнения математической физики
\publaddr М.
\publ Наука
\yr 1978
\totalpages 392 
\zmath{https://zbmath.org/?q=an:0384.35001}

\RBibitem{zayn:7}
\by Сабитов~К.~Б.
\book Уравнения математической физики
\publaddr М.
\publ Физматлит
\yr 2013
\totalpages 352 



\RBibitem{zayn:8}
\by Сабитов~К.~Б.
\paper Об одной краевой задаче для уравнения смешанного типа третьего порядка
\jour Докл. РАН
\yr 2009
\vol 427
\issue 5
\pages 593--596
\zmath{https://zbmath.org/?q=an:1190.35167}
\elib{http://elibrary.ru/item.asp?id=12803887}





\RBibitem{zayn:9}
\by Зайнуллов~А.~Р.
\paper Обратные задачи для уравнения теплопроводности
\jour Вестн. СамГУ. Естественнонаучн. сер.
\yr 2015
\issue 6(128)
\pages 62--75


\RBibitem{zayn:10}
\by Тихонов~А.~Н., Арсенин~В.~Я.
\book Методы решения некорректных задач 
\publaddr   М.
\publ Наука
\yr 1979
\totalpages 285 

\RBibitem{zayn:11}
\by Лавретьев~М.~М., Романов~В.~Г., Шишатский~С.~П.
\book Некорректные задачи математической физики и анализа
\publ Наука
\publaddr М.
\yr 1980
\totalpages 286 
\zmath{https://zbmath.org/?q=an:0476.35001}

\RBibitem{zayn:12}
\by Иванов~В.~К., Васин~В.~В., Танана~В.~П.
\book Теория линейных некорректных задач и её приложения
\publaddr М.
\publ Наука
\yr 1978
\totalpages 206 
\zmath{https://zbmath.org/?q=an:03770864}


\end{thebibliography}

\RuDate

\newpage

\makeentitle



\thanks{\textbf{Acknowledgments.} 
This work was  supported by the Russian Foundation for Basic Research  (project no.~14--01--97003-r\_povolzh'e\_a).}



\thanks{{\bf ORCID}

{\bf Artur R. Zaynullov:}    \url{http://orcid.org/0000-0003-2657-4101}



}




\EngRef

\Russian

\begin{thebibliography}{99}

\Bibitem{zayn:1:eng}
\by Denisov~A.~M.
\book Elements of the theory of inverse problems
\serial Inverse and Ill-Posed Problems Series
\publaddr Utrecht
\publ VSP
\yr 1999.
%\totalpages 
iv+272~pp
\zmath{https://zbmath.org/?q=an:0940.35003}
\crossref{10.1515/9783110943252}

\Bibitem{zayn:2:eng}
\by Tikhonov~A.~N.; Samarskii~A.~A.
\book Equations of mathematical physics
\serial International Series of Monographs on Pure and Applied Mathematics
\vol 39
\publaddr Oxford etc.
\publ Pergamon Press
\yr 1963
\totalpages xvi+765 
\zmath{https://zbmath.org/?q=an:0111.29008}


\Bibitem{zayn:3:eng}
\lang In Russian
\by  Sobolev~S.~L.
\book Uravneniia matematicheskoi fiziki \rm [Equations of mathematical physics]
\publaddr Moscow
\publ Nauka
\yr 1992
\totalpages 432


\Bibitem{zayn:4:eng}
\lang In Russian
\by Vladimirov~V.~S., Zharinov~V.~V.
\book Uravneniia matematicheskoi fiziki \rm [Equations of mathematical physics]
\publaddr M.
\publ Fizmatlit
\yr 2000
\totalpages 398

\Bibitem{zayn:5:eng}
\by Koshlyakov~N.~S., Smirnov~M.~M., Gliner~E.~B.
\book Differential equations of mathematical physics
\publaddr Amsterdam
\publ North-Holland Publishing Company
\totalpages xvi+701 
\yr 1964
\zmath{https://zbmath.org/?q=an:0115.30701}

\Bibitem{zayn:6:eng}
\lang In Russian
\by Godunov~S.~K.
\book Uravneniia matematicheskoi fiziki \rm [Equations of mathematical physics]
\publaddr Moscow
\publ Nauka
\yr 1978
\totalpages 392 


\Bibitem{zayn:7:eng}
\lang In Russian
\by Sabitov~K.~B.
\book Uravneniia matematicheskoi fiziki \rm [Equations of mathematical physics]
\publaddr Moscow
\publ Fizmatlit
\yr 2013
\totalpages 352 


\Bibitem{zayn:8:eng}
\by Sabitov~K.~B.
\paper A boundary value problem for a third-order equation of mixed type
\jour Dokl. Math.
\yr 2009
\vol 80
\issue 1
\pages 565--568
\crossref{10.1134/S1064562409040292}
\mathscinet{http://www.ams.org/mathscinet-getitem?mr=2590023}
\zmath{https://zbmath.org/?q=an:1190.35167}
\isi{http://gateway.isiknowledge.com/gateway/Gateway.cgi?GWVersion=2&SrcApp=PARTNER_APP&SrcAuth=LinksAMR&DestLinkType=FullRecord&DestApp=ALL_WOS&KeyUT=000269308500029}
\elib{http://elibrary.ru/item.asp?id=15298865}


\Bibitem{zayn:9:eng}
\lang In Russian
\by Zaynullov~A.~R.
\paper Inverse problems for the heat equation
\jour Vestnik SamGU. Estestvenno-Nauchnaya Ser.
\yr 2015
\issue 6(128)
\pages 62--75


\Bibitem{zayn:10:eng}
\by Tikhonov~A.~N., Arsenin~V.~Ya.
\book Solutions of ill-posed problems
\zmath{https://zbmath.org/?q=an:0354.65028}
\serial Scripta Series in Mathematics
\publaddr New York etc.
\publ John Wiley \& Sons
\totalpages xiii+258 
\yr 1977

\Bibitem{zayn:11:eng}
\by Lavrent'ev~M.~M., Romanov~V.~G., Shishatskij~S.~P.
\book Ill-posed problems of mathematical physics and analysis
\zmath{https://zbmath.org/?q=an:0593.35003}
\serial Translations of Mathematical Monographs
\vol  64
\publaddr Providence, R.I.
\publ American Mathematical Society 
\totalpages vi+290 
\yr 1986

\Bibitem{zayn:12:eng}
\by Ivanov~V.~K., Vasin~V.~V., Tanana~V.~P.
\book Theory of linear ill-posed problems and its applications
\serial Inverse and Ill-Posed Problems Series
\publaddr Utrecht
\publ VSP 
\yr 2002
%\totalpages 
xiii+281~pp 
\zmath{https://zbmath.org/?q=an:1037.65056}
\crossref{10.1515/9783110944822}


\end{thebibliography}

\RuDate


\newpage
\Russian

%
%\end{document}
